\section{Sujet n\textsuperscript{o}\,12}
\noindent\begin{minipage}{0.5\linewidth}\footnotesize Office du Baccalauréat du Cameroun\end{minipage}%
\hfill\begin{minipage}{0.45\linewidth}\footnotesize\raggedleft Examen : \textbf{Baccalauréat} · Série : \textbf{C/E}\\
Épreuve : \textbf{Mathématiques} · Durée : \textbf{4\,h} · Coef. : \textbf{7}\end{minipage}
\par\smallskip{\color{onbo}\rule{\linewidth}{1pt}}

\subsection*{Partie A — Évaluation des ressources \hfill (15 points)}

\noindent\textbf{Exercice 1.} \hfill\textit{(3 points)}\\
Dans une huilerie d'Édéa, un lot de contrôle est constitué de $6$ régimes de palme
dont $2$ sont défectueux (mal mûris), les régimes étant indiscernables au toucher.
Un technicien prélève simultanément $3$ régimes de ce lot. On note $X$ le nombre de
régimes défectueux figurant parmi les $3$ régimes prélevés.
\begin{enumerate}[label=\arabic*)]
\item Déterminer la loi de probabilité de $X$. \hfill\textit{(1,5 pt)}
\item Calculer l'espérance mathématique $E(X)$ et interpréter ce résultat. \hfill\textit{(0,75 pt)}
\item Calculer la variance $V(X)$ et l'écart-type $\sigma(X)$. \hfill\textit{(0,75 pt)}
\end{enumerate}

\smallskip
\noindent\textbf{Exercice 2.} \hfill\textit{(3 points)}\\
La quantité d'huile brute (en tonnes) stockée dans une cuve tampon le jour $n$ est
modélisée par la suite $(u_n)$ définie par $u_0=200$ et, pour tout $n\in\N$,
$u_{n+1}=\dfrac45\,u_n+30$.
\begin{enumerate}[label=\arabic*)]
\item On pose $v_n=u_n-150$. Démontrer que $(v_n)$ est géométrique ; préciser sa raison et son premier terme. \hfill\textit{(1 pt)}
\item Exprimer $v_n$ puis $u_n$ en fonction de $n$. \hfill\textit{(1 pt)}
\item Étudier la convergence de $(u_n)$ et interpréter la limite obtenue. \hfill\textit{(1 pt)}
\end{enumerate}

\smallskip
\noindent\textbf{Exercice 3.} \hfill\textit{(4 points)}\\
Soit $f$ la fonction définie sur $]0\,;\,+\infty[$ par $f(x)=\dfrac{\ln x}{x}$, de courbe
représentative $(\mathcal{C})$ dans un repère orthonormé.
\begin{enumerate}[label=\arabic*)]
\item Déterminer les limites de $f$ en $0^+$ et en $+\infty$ ; interpréter graphiquement chacune. \hfill\textit{(1 pt)}
\item Calculer $f'(x)$, dresser le tableau de variations de $f$ et préciser le maximum de $f$. \hfill\textit{(1,5 pt)}
\item À l'aide d'une intégration par parties, calculer $\displaystyle\int_1^{\mathrm{e}}\ln x\,\dd x$. \hfill\textit{(0,75 pt)}
\item En déduire, en unités d'aire, l'aire $\mathcal{A}$ du domaine limité par $(\mathcal{C})$,
l'axe des abscisses et les droites d'équations $x=1$ et $x=\mathrm{e}$. \hfill\textit{(0,75 pt)}
\end{enumerate}

\smallskip
\noindent\textbf{Exercice 4.} \hfill\textit{(5 points)}\\
Le plan complexe est muni d'un repère orthonormé direct $(O\,;\,\vec u,\vec v)$.
\begin{enumerate}[label=\arabic*)]
\item Résoudre dans $\C$ l'équation $z^2-2\sqrt3\,z+4=0$ ; on note $z_1$ la solution de
partie imaginaire positive et $z_2$ l'autre. \hfill\textit{(1 pt)}
\item Écrire $z_1$ et $z_2$ sous forme trigonométrique, puis calculer $z_1^{\,6}$. \hfill\textit{(1,25 pt)}
\item Les points $A$, $B$, $C$ ont pour affixes respectives $z_1$, $z_2$ et $2\sqrt3$.
Calculer $AB$, $AC$ et $BC$, puis préciser la nature du triangle $ABC$. \hfill\textit{(1,25 pt)}
\item Soit $S$ la similitude directe de centre $O$, de rapport $2$ et d'angle $\frac\pi3$.
Donner l'écriture complexe de $S$, puis déterminer l'affixe de l'image $A'$ du point $A$ par $S$. \hfill\textit{(1,5 pt)}
\end{enumerate}

\subsection*{Partie B — Évaluation des compétences \hfill (5 points)}

\noindent\textit{Monsieur Ekani dirige une petite unité d'extraction d'huile de palme à
Édéa et te consulte, en tant qu'élève de Terminale~C, pour trois décisions. Il veut
d'abord fabriquer, en tôle, une cuve de stockage ayant la forme d'un pavé droit à
\emph{base carrée}, \emph{ouverte sur le dessus} ; il dispose d'exactement
\SI{12}{\metre\squared} de tôle et souhaite que la cuve ait le plus grand volume
possible. Par ailleurs, sa production d'huile, de \SI{200}{\tonne} cette année,
augmente de \SI{8}{\percent} chaque année. Enfin, sur sa chaîne, \SI{5}{\percent} des
bidons sont mal scellés ; un grossiste en prélève $8$ au hasard à la livraison.}

\begin{enumerate}[label=\textbf{Tâche \arabic*.},leftmargin=2.4em]
\item Déterminer les dimensions (côté de la base et hauteur) de la cuve de \emph{volume maximal}
qu'il peut fabriquer, ainsi que ce volume. \hfill\textit{(1,5 pt)}
\item Déterminer l'année à partir de laquelle sa production annuelle dépassera \SI{300}{\tonne}. \hfill\textit{(1,5 pt)}
\item Déterminer la probabilité qu'\emph{au moins un} des $8$ bidons prélevés soit mal scellé. \hfill\textit{(1,5 pt)}
\end{enumerate}
\par\smallskip{\footnotesize\itshape Présentation générale : 0,5 point.}

% ════════════════════════════════════════════════════
\corrige{du sujet 12}

\noindent\textbf{Partie A — Exercice 1.}
1) Il y a $\binom63=20$ tirages équiprobables. $X$ prend les valeurs $0,1,2$ et
$P(X{=}k)=\dfrac{\binom2k\binom4{3-k}}{20}$ :
$P(0)=\frac{\binom40}{20}\cdot\binom20=\frac{4}{20}=\frac15$,
$P(1)=\frac{\binom21\binom42}{20}=\frac{12}{20}=\frac35$,
$P(2)=\frac{\binom22\binom41}{20}=\frac{4}{20}=\frac15$ (somme $=1$).
2) $E(X)=0\cdot\frac15+1\cdot\frac35+2\cdot\frac15=\frac35+\frac25=1$ : en moyenne, un
prélèvement de $3$ régimes contient $1$ régime défectueux.
3) $E(X^2)=0+1\cdot\frac35+4\cdot\frac15=\frac35+\frac45=\frac75$ ; donc
$V(X)=E(X^2)-E(X)^2=\frac75-1=\frac25$ et $\sigma(X)=\sqrt{\frac25}=\frac{\sqrt{10}}{5}\approx0{,}63$.

\noindent\textbf{Exercice 2.}
1) $v_{n+1}=u_{n+1}-150=\frac45u_n+30-150=\frac45u_n-120=\frac45(u_n-150)=\frac45v_n$.
Donc $(v_n)$ est géométrique de raison $q=\frac45$ et de premier terme $v_0=u_0-150=50$.
2) $v_n=50\left(\frac45\right)^{\!n}$, d'où $u_n=v_n+150=150+50\left(\frac45\right)^{\!n}$.
3) Comme $0<\frac45<1$, $\left(\frac45\right)^{n}\to0$, donc $u_n\to150$. Le stock se
stabilise au fil du temps autour de \SI{150}{\tonne}.

\noindent\textbf{Exercice 3.}
1) En $0^+$ : $\ln x\to-\infty$ et $\frac1x\to+\infty$, donc $f(x)\to-\infty$ : l'axe des
ordonnées ($x=0$) est asymptote verticale. En $+\infty$ : par croissances comparées
$\frac{\ln x}{x}\to0^+$, donc la droite $y=0$ est asymptote horizontale.
2) $f'(x)=\dfrac{\frac1x\cdot x-\ln x\cdot1}{x^2}=\dfrac{1-\ln x}{x^2}$. Sur $]0;+\infty[$,
$f'(x)>0\iff\ln x<1\iff x<\mathrm{e}$. Donc $f$ croît sur $]0;\mathrm{e}]$, décroît sur
$[\mathrm{e};+\infty[$, avec un maximum $f(\mathrm{e})=\frac{\ln\mathrm{e}}{\mathrm{e}}=\frac1{\mathrm{e}}$.
3) IPP avec $u=\ln x$, $v'=1$, donc $u'=\frac1x$, $v=x$ :
$\int_1^{\mathrm{e}}\ln x\,\dd x=\big[x\ln x\big]_1^{\mathrm{e}}-\int_1^{\mathrm{e}}1\,\dd x
=(\mathrm{e}-0)-\big[x\big]_1^{\mathrm{e}}=\mathrm{e}-(\mathrm{e}-1)=1$.
4) Sur $[1;\mathrm{e}]$, $\ln x\ge0$ donc $f(x)\ge0$ : l'aire vaut
$\mathcal{A}=\int_1^{\mathrm{e}}\frac{\ln x}{x}\,\dd x=\Big[\frac{(\ln x)^2}{2}\Big]_1^{\mathrm{e}}
=\frac{1}{2}-0=\frac12$ unité d'aire.

\noindent\textbf{Exercice 4.}
1) $\Delta=(2\sqrt3)^2-4\times4=12-16=-4=(2i)^2$. Les racines sont
$z=\frac{2\sqrt3\pm2i}{2}=\sqrt3\pm i$, d'où $z_1=\sqrt3+i$ et $z_2=\sqrt3-i$.
2) $|z_1|=\sqrt{3+1}=2$ ; $z_1=2\!\left(\frac{\sqrt3}{2}+\frac i2\right)=2\,\mathrm{e}^{i\pi/6}$
et $z_2=\bar{z_1}=2\,\mathrm{e}^{-i\pi/6}$. Donc
$z_1^{\,6}=2^6\,\mathrm{e}^{i\pi}=64\times(-1)=-64$.
3) $AB=|z_2-z_1|=|-2i|=2$ ; $AC=|2\sqrt3-z_1|=|\sqrt3-i|=2$ ;
$BC=|2\sqrt3-z_2|=|\sqrt3+i|=2$. Comme $AB=AC=BC=2$, le triangle $ABC$ est \textbf{équilatéral}.
4) $S(z)=2\,\mathrm{e}^{i\pi/3}z=2\!\left(\frac12+\frac{\sqrt3}{2}i\right)z=(1+i\sqrt3)\,z$.
Affixe de $A'$ : $(1+i\sqrt3)(\sqrt3+i)=\sqrt3+i+3i+i^2\sqrt3=\sqrt3-\sqrt3+(1+3)i=4i$.

\smallskip
\noindent\textbf{Partie B — Situation.}
\emph{Tâche 1.} Notons $x>0$ le côté de la base (en m) et $h$ la hauteur. La tôle (base
$+$ $4$ faces latérales) vérifie $x^2+4xh=12$, d'où $h=\dfrac{12-x^2}{4x}$ (avec $0<x<2\sqrt3$).
Le volume est $V(x)=x^2h=\dfrac{x(12-x^2)}{4}=\dfrac{12x-x^3}{4}$. Alors
$V'(x)=\dfrac{12-3x^2}{4}=\dfrac{3(4-x^2)}{4}$, qui s'annule en $x=2$ (positif), avec
$V'>0$ sur $]0;2[$ et $V'<0$ sur $]2;2\sqrt3[$ : maximum en $x=2$. On obtient alors
$h=\frac{12-4}{8}=1$ et $V(2)=\frac{24-8}{4}=4$. La cuve optimale a une base de
\SI{2}{\metre} de côté, une hauteur de \SI{1}{\metre}, pour un volume maximal de \SI{4}{\metre\cubed}.

\emph{Tâche 2.} La production de l'année $n$ (avec $n=0$ pour cette année) est
$p_n=200\times1{,}08^{\,n}$. On résout $p_n>300$, soit $1{,}08^{\,n}>\frac32$, d'où
$n>\dfrac{\ln(3/2)}{\ln1{,}08}\approx5{,}27$ : c'est donc à partir de la \textbf{6\textsuperscript{e} année}
suivant l'année actuelle ($p_6\approx\SI{317}{\tonne}$, alors que $p_5\approx\SI{294}{\tonne}$).

\emph{Tâche 3.} Chaque bidon est mal scellé avec la probabilité $0{,}05$,
indépendamment des autres ; le nombre $Y$ de bidons mal scellés parmi $8$ suit une loi
binomiale $\mathcal{B}(8\,;\,0{,}05)$. Donc
$P(Y\ge1)=1-P(Y{=}0)=1-(0{,}95)^8\approx1-0{,}663=\mathbf{0{,}337}$, soit environ \SI{33,7}{\percent}.
